\section{Plan de pruebas}

Las pruebas unitarias están en el mismo archivo que cada módulo
(\texttt{actores}, \texttt{heartbeat}, \texttt{bully},
\texttt{dos\_fases}, \texttt{petri}) y se ejecutan con
\texttt{cargo test --lib}.

En actores se verifica que muchos débitos concurrentes no dejen saldo
negativo, que un débito sin fondos no modifique la cuenta y que una sola
entrada no se asigne dos veces. En heartbeats, que la víctima pase por
sospechoso antes de caído, que un nodo sano no genere eventos y que una
falsa alarma se recupere. En Bully, que gane el vivo de mayor ID con el
líder caído, con dos mayores caídos y con todos vivos. En 2PC se cubren
la compra exitosa, los \texttt{No} locales, el timeout de un
participante y la caída del coordinador después del \texttt{Prepare},
que deja a ambos bloqueados. En Petri se comprueba que el camino
exitoso conserve el invariante y la 1-acotación.

El ejecutable (\texttt{parte1} \ldots{} \texttt{parte4}) reproduce esos
escenarios en consola. Si el coordinador cae tras el \texttt{Yes}, el
token del participante queda en \texttt{P\_Prepared} sin transición
habilitada: es el mismo estado \texttt{bloqueado} de
\texttt{dos\_fases}.
